% !TEX program = xelatex
% GENERATED by tools/build.py from data/records.csv and templates/cv_short.tex.
% One A4 page. The build stops if this runs onto a second page.
\documentclass[10pt, a4paper]{cvClass}
\geometry{left=1.6cm, top=1.0cm, right=1.6cm, bottom=0.8cm, footskip=.3cm}
% ---------------------------------------------------------------------------
% Shared layout macros for cv_long.tex, cv_short.tex and publication_list.tex
% (input by each template; lives in templates/, copied next to the .tex)
% ---------------------------------------------------------------------------
\colorlet{awesome}{awesome-red}          % the CV orange-red (#DC3522)
\definecolor{accent2}{HTML}{E3B23C}      % muted yellow, used sparingly
\setlength{\parskip}{0pt}
\urlstyle{same}

% A hanging list: year (or label) on the left, text on the right
\newlength{\cvlabelwidth}\setlength{\cvlabelwidth}{2.45cm}
\newlength{\cvitemsep}\setlength{\cvitemsep}{1.4mm}
\newenvironment{cvlist}{%
  \vspace{1.2mm}%
  \begin{list}{}{%
    \setlength{\leftmargin}{\cvlabelwidth}\setlength{\labelwidth}{\dimexpr\cvlabelwidth-3mm\relax}%
    \setlength{\labelsep}{3mm}\setlength{\itemsep}{\cvitemsep}%
    \setlength{\parsep}{0pt}\setlength{\topsep}{0pt}\setlength{\partopsep}{0pt}}%
  \fontsize{9pt}{11.5pt}\bodyfontlight\color{text}\raggedright}%
  {\end{list}\vspace{1mm}}
\newcommand{\cvitem}[2]{\item[{\fontsize{8.5pt}{10pt}\bodyfont\color{graytext}#1\hfill}]#2}
\newcommand{\skilllabel}[1]{{\bodyfont\color{darktext}#1}}
\renewcommand{\textbf}[1]{{\bodyfont\bfseries\color{darktext}#1}}

% tighter spacing for the one-page version
\renewcommand{\acvSectionTopSkip}{1.2mm}
\renewcommand{\acvSectionContentTopSkip}{1.5mm}
\renewcommand{\acvHeaderAfterSocialSkip}{1mm}
\setlength{\cvitemsep}{0.5mm}
\setlength{\cvlabelwidth}{2.35cm}
\renewcommand*{\sectionstyle}[1]{{\fontsize{12pt}{1em}\bodyfont\bfseries\color{awesome}#1}}

\photo[circle,edge,left]{Tobias_Staal_square}
\name{Tobias}{Stål}
\position{Geoscientist and engineer}
\address{Tasmania, Australia}
\mobile{+61 4 4749 4089}
\email{tobias.staal@utas.edu.au}
\homepage{tobbetripitaka.github.io/Tobias\_Staal}
\github{TobbeTripitaka}
\makecvfooter{}{\color{lighttext}Full CV and publication list: tobbetripitaka.github.io/Tobias\_Staal}{}

\begin{document}
\makecvheader

\vspace{3mm}
{\fontsize{9pt}{11.5pt}\bodyfontlight\color{text}Geophysicist working on Antarctica's ice-covered interior. I integrate geology, geophysics and statistical modelling to understand interactions between the ice sheet and the solid Earth, with a focus on geothermal heat, subglacial geology and instrumentation. Background in geology and geomorphology from the Arctic and Europe, and an earlier professional career in sound, lighting and electrical engineering.\par}

\cvsection{Academic positions}
\begin{cvlist}
\cvitem{2022 – Present}{\textbf{Research Associate in Computational Physics (Geophysics)}, ACEAS, University of Tasmania, Tasmania, Australia}
\cvitem{2021 – 2022}{\textbf{Research Fellow in Antarctic Seismology}, University of Tasmania, Tasmania, Australia}
\end{cvlist}

\cvsection{Selected non-academic work experience}
\begin{cvlist}
\cvitem{2021 – Present}{\textbf{GRIT Facility Manager – Instrumentation}, University of Tasmania, Tasmania, Australia. Instrumentation management and Antarctic fieldwork support}
\cvitem{2010 – 2015}{\textbf{Field Leader and Geophysicist}, COWI, Denmark. MEP, GPR and TEM fieldwork and survey design}
\cvitem{1997 – 2016}{\textbf{Sound engineer, sound, lighting and video designer, technical producer}, Swedish Radio, Royal Dramatic Theatre, Circus Cirkör, National Theatre of Ghana, STÜ and others, Sweden, Denmark, Estonia, Ghana, Palestine}
\end{cvlist}

\cvsection{Education}
\begin{cvlist}
\cvitem{2021}{\textbf{PhD, Geophysics, Geology and Computational Methods}, University of Tasmania, Australia}
\cvitem{2015}{\textbf{MSc, Geophysics}, University of Copenhagen, Denmark}
\cvitem{2011}{\textbf{BSc, Geology}, University of Copenhagen, Denmark}
\cvitem{1998}{\textbf{AD, Sound Engineering}, Piteå Music Conservatory (Luleå University of Technology), Piteå, Sweden}
\end{cvlist}

\cvsection{Selected publications}
\begin{cvlist}
\cvitem{2025}{NJ Abram, A Purich, …, \textbf{T Stål} et al. — Emerging evidence of abrupt changes in the Antarctic environment. \textit{Nature}.}
\cvitem{2024}{\textbf{T Stål}, JA Halpin, JW Goodge et al. — Geology Matters for Antarctic Geothermal Heat. \textit{Geophysical Research Letters}.}
\cvitem{2022}{\textbf{T Stål}, AM Reading, S Fuchs et al. — Properties and biases of the global heat flow compilation. \textit{Frontiers in Earth Science}.}
\cvitem{2022}{AM Reading, \textbf{T Stål}, JA Halpin et al. — Antarctic geothermal heat flow and its implications for tectonics and ice sheets. \textit{Nature Reviews Earth \& Environment}.}
\cvitem{2021}{\textbf{T Stål}, AM Reading, JA Halpin et al. — Antarctic Geothermal Heat Flow Model: Aq1. \textit{Geochemistry, Geophysics, Geosystems}.}
\cvitem{2020}{\textbf{T Stål}, AM Reading, JA Halpin et al. — The Antarctic Crust and Upper Mantle: A Flexible 3D Model and Software Framework for Interdisciplinary Research. \textit{Frontiers in Earth Science}.}
\cvitem{2019}{\textbf{T Stål}, AM Reading, JA Halpin et al. — A Multivariate Approach for Mapping Lithospheric Domain Boundaries in East Antarctica. \textit{Geophysical Research Letters}.}
\end{cvlist}

\cvsection{Datasets}
\begin{cvlist}
\cvitem{2021}{\textbf{T Stål}, AM Reading, S Kupis et al. — Casey-Wilkins Adaptable Array [Dataset]. \textit{International Federation of Digital Seismograph Networks}.}
\end{cvlist}

\cvsection{Awards and recognition}
\begin{cvlist}
\cvitem{2023}{Royal Society of Tasmania Doctoral Award}
\cvitem{2021}{Wiley Top Cited Article}
\cvitem{2020}{PhD Award for Outstanding Performance during Postgraduate Studies}
\end{cvlist}

\cvsection{Memberships and service}
\begin{cvlist}
\cvitem{2023 – Present}{Member, International Heat Flow Commission (IHFC), IUGG}
\cvitem{2021 – Present}{Co-chair, SCAR INSTANT Heat Flow Subcommittee}
\end{cvlist}

\cvsection{Art collaborations (as a scientist)}
\begin{cvlist}
\cvitem{2025}{T Stål, AM Reading, K Verell — \textit{Seismic Dance Party}, Beaker Street (Hobart). Denman Glacier seismic data as an audio-visual dancefloor experience.}
\cvitem{2025}{L Bleach — \textit{An Infrasonorous Archipelago}, Walk\&Talk Biennial "Gestures of Abundance" (Azores). Seismo-volcanic infrasound as tactile and sonic works.}
\end{cvlist}

\cvsection{Other skills and experience}
\begin{cvlist}
\cvitem{\skilllabel{Coding}}{Python (scikit-learn, TensorFlow, dask, xarray, geopandas), GIS (QGIS, ArcGIS, GDAL), LaTeX/TikZ, Bash, git, SCons, SQL, Seismic Unix, Madagascar, Blender}
\cvitem{\skilllabel{Certificates}}{Antarctic trip leader and terrain-vehicle driver (AAD, Casey 2023); Wilderness First Aid; Coastal Skipper; PADI Advanced Open Water}
\cvitem{\skilllabel{Languages}}{English (C2), Swedish (native), Danish (fluent), French (intermediate), Russian and Asante Twi (functional)}
\end{cvlist}


\end{document}
